\documentclass[10pt]{article}
\usepackage[T1]{fontenc}
\usepackage[utf8]{inputenc}
\usepackage[french]{babel}
\usepackage{amsmath,amssymb,mathrsfs,mathpazo}
\usepackage{geometry}
\geometry{a4paper,margin=2cm}
\usepackage{graphicx,array,tabularx,tabularray,colortbl,multirow,multicol,booktabs,enumitem,pifont,ulem,fancybox,diagbox,cancel}
\usepackage{pstricks,pstricks-add,pst-plot,pst-tree,pst-node,pst-text,pst-func,pst-math,pst-bspline,pst-3dplot,pst-eucl,pst-solides3d}
\let\pscalculate\undefined
\usepackage{pst-bezier}
\usepackage{tikz,pgfplots,tkz-tab}
\usetikzlibrary{arrows,arrows.meta,patterns,calc,positioning,trees,decorations.markings,shapes.geometric,shapes.misc,shapes.arrows,intersections,decorations.pathmorphing,decorations.pathreplacing,backgrounds,shadows,babel}
\pgfplotsset{compat=1.18}
\usepgfplotslibrary{fillbetween}
\usepackage[np]{numprint}
\usepackage{listings,makecell,pas-tableur,twoopt}
\let\barre\undefined
\usepackage[pstricks]{tablvar}
\definecolor{darkgreen}{RGB}{0,100,0}
\providecommand{\red}{\color{red}}
\providecommand{\blue}{\color{blue}}
\providecommand{\green}{\color{green}}
\providecommand{\black}{\color{black}}
\providecommand{\white}{\color{white}}
\providecommand{\gray}{\color{gray}}
\providecommand{\cyan}{\color{cyan}}
\providecommand{\magenta}{\color{magenta}}
\providecommand{\yellow}{\color{yellow}}
\providecommand{\violet}{\color{violet}}
\providecommand{\orange}{\color{orange}}
\providecommand{\euro}{\texteuro}
\providecommand{\EUR}{\texteuro}
\providecommand{\ieme}{\textsuperscript{e}}
\providecommand{\degres}{\ensuremath{^\circ}}
\providecommand{\degree}{\ensuremath{^\circ}}
\providecommand{\vv}[1]{\overrightarrow{#1}}
\providecommand{\Coord}[1]{#1}
\providecommand{\decosix}{}
\providecommand{\starredbullet}{\ensuremath{\bullet}}
\pagestyle{empty}
\setlength{\parindent}{0pt}
\usepackage{adjustbox}\DeclareUnicodeCharacter{2236}{:}\DeclareUnicodeCharacter{2190}{\ensuremath{\leftarrow}}
\providecommand{\R}{\mathbb{R}}
\providecommand{\N}{\mathbb{N}}
\providecommand{\Z}{\mathbb{Z}}
\providecommand{\Q}{\mathbb{Q}}
\providecommand{\D}{\mathbb{D}}
\providecommand{\C}{\mathbb{C}}
\providecommand{\e}{{\mathrm{e}}}
\providecommand{\vect}[1]{\overrightarrow{#1}}
\providecommand{\vectt}[1]{\overrightarrow{\mathrm{#1}}}
\providecommand{\vvt}[1]{\overrightarrow{\mathrm{#1}}}
\providecommand{\np}[2][]{{#2~#1}}
\providecommand{\vv}[1]{{\overrightarrow{#1}}}
\renewcommand{\d}{{\mathrm{d}}}
\providecommand{\degres}{{^\circ}}
\providecommand{\degre}{{^\circ}}
\providecommand{\blue}{\color{blue}}
\providecommand{\red}{\color{red}}
\providecommand{\enskip}{\quad}
\providecommand{\llbracket}{\lbrack\!\lbrack}
\providecommand{\rrbracket}{\rbrack\!\rbrack}
\providecommand{\euro}{{\text{€}}}
\providecommand{\no}{{\text{n°}}}
\providecommand{\interval}[2]{{[#1;#2]}}
\providecommand{\intervFO}[2]{{[#1;#2[}}
\providecommand{\dotfill}{\ldots}

\providecommand{\ds}{\displaystyle}
\providecommand{\textding}[1]{#1}
\providecommand{\Oij}{\(\mathrm{O};\vec{\imath},\vec{\jmath}\)}
\providecommand{\Oijk}{\(\mathrm{O};\vec{\imath},\vec{\jmath},\vec{k}\)}
\providecommand{\Ouv}{\(\mathrm{O};\vec{u},\vec{v}\)}
\providecommand{\intervFF}[2]{\ensuremath{[#1;#2]}}
\renewcommand{\euro}{\ifmmode\text{\texteuro}\else\texteuro\fi}
\DeclareUnicodeCharacter{00B0}{\ifmmode^{\circ}\else\textdegree\fi}
\DeclareUnicodeCharacter{20AC}{\ifmmode\text{\texteuro}\else\texteuro\fi}
% Énoncés : Éducation nationale. Transcription et corrigés : APMEP, Denis Vergès et contributeurs. Découpage et classement : CoopMaths, Sébastien Lozano, Éric Elter et contributeurs. Adaptation de la présentation : Maths Explorer.
% CC BY-SA : https://coopmaths.fr/www/about/
% https://coopmaths.fr/alea/static/e3c/2020/tex/e3c_2020_00_sujet42_3.tex
\begin{document}
\noindent\textit{APMEP / CoopMaths. CC BY-SA. Présentation adaptée par Maths Explorer.}\par\medskip
\medskip

\textbf{Partie A}

\medskip

Étudier sur $\R$ le signe de $P(x)= - 10x^2 - 40x + 120$.

\medskip

\textbf{Partie B}

\medskip

On se place dans un plan muni d'un  repère orthonormé. La courbe $H$ représentée sur le graphique ci -dessous est l'ensemble des points de l'hyperbole d'équation:

\[y = \dfrac{10x + 4}{x + 2}\]

avec $x$ appartenant à l'intervalle [0~;~8].

Pour toute abscisse $x$ dans l'intervalle [0~;~8], on construit le rectangle ABDE comme indiqué sur la figure. On donne les informations suivantes :

\medskip

\parbox{0.35\linewidth}{
\setlength\parindent{1cm}
\begin{itemize}[label=\textbullet]
\item A et B sont sur l'axe des abscisses;
\item A est d'abscisse $x$;
\item B et D ont pour abscisse 8 ;
\item E appartient à la courbe H ;
\item D et E ont la même ordonnée.
\end{itemize}
\setlength\parindent{0cm}
}\hfill
\parbox{0.63\linewidth}{\psset{unit=0.7cm}
\begin{pspicture}(-1,-1)(9,9)
\psframe[fillstyle=solid,fillcolor=lightgray](1,0)(8,4.66667)
\psgrid[gridlabels=0pt,subgriddiv=1](0,0)(9,9)
\psaxes[linewidth=1.25pt,labelFontSize=\scriptstyle]{->}(0,0)(0,0)(9,9)
\psplot[plotpoints=2000,linewidth=1.25pt,linecolor=red]{0}{8}{10 x mul 4 add x 2 add div}
\uput[dl](1,0){\footnotesize A}\uput[dr](8,0){\footnotesize B}
\uput[ur](8,4.66667){\footnotesize D}\uput[ul](1,4.66667){\footnotesize E}
\uput[u](3.2,7){\red $H$}
\end{pspicture}}

\medskip

L'objectif de ce problème est de déterminer la ou les valeurs éventuelles $x$ de l'intervalle [0~;~8] correspondant à un rectangle ABDE d'aire maximale.

\medskip

\begin{enumerate}
\item Déterminer l'aire du rectangle ABDE lorsque $x = 0$.
\item Déterminer l'aire du rectangle ABDE lorsque $x =  4$.
\end{enumerate}

On définit la fonction $f$ qui à tout réel $x$ de [0~;~8], associe l'aire du rectangle ABDE. On admet que:

\[f(x) = \dfrac{-10x^2 + 76x + 32}{x + 2}.\]

\begin{enumerate}[resume]
\item Répondre au problème posé.
\end{enumerate}

\bigskip

\end{document}
